\section{Méthodes d'extraction de caractéristiques}
\label{sec:features}

Quatre familles de descripteurs complémentaires sont extraites de chaque image
prétraitée. Chacune capture un aspect distinct de l'apparence tumorale.

\subsection{A --- Caractéristiques statistiques (distribution d'intensité)}

Onze statistiques du premier ordre sont calculées à partir de l'histogramme d'intensités
des pixels : moyenne, écart-type, variance, entropie de Shannon, asymétrie, kurtosis
et cinq percentiles (p10, p25, p50, p75, p90). Elles fournissent un résumé compact
de la distribution globale d'intensité, sans structure spatiale.

\subsection{B --- Motifs binaires locaux (LBP)}

Le LBP encode la texture locale en comparant chaque pixel à son voisinage circulaire
de $P$ voisins au rayon $R$. La méthode d'encodage détermine le regroupement en histogramme.
Le LBP capture les micro-motifs de texture liés à l'hétérogénéité tissulaire tumorale.

\subsection{C --- Transformée en ondelettes discrète (DWT)}

Une DWT 2D à un niveau décompose l'image en quatre sous-bandes (LL, LH, HL, HH)
représentant respectivement la structure grossière et les détails haute fréquence
horizontaux, verticaux et diagonaux. Des statistiques sont extraites de chaque
sous-bande, fournissant une information multi-échelle dans le domaine fréquentiel.

\subsection{D --- Histogrammes de gradients orientés (HOG)}

Le HOG capture la distribution des magnitudes et orientations des gradients locaux
en divisant l'image en cellules et en calculant des histogrammes d'orientation,
puis en normalisant sur des blocs chevauchants. Le HOG est bien établi pour la
discrimination basée sur la forme et les contours.

\bigskip
La Figure~\ref{fig:descriptor_demo} illustre les quatre types de descripteurs
appliqués à une même coupe IRM.

\begin{figure}[H]
  \centering
  \includegraphics[width=\linewidth]{descriptor_demo.png}
  \caption{Descripteurs artisanaux calculés à partir d'une coupe IRM prétraitée
           (méningiome). Haut : coupe prétraitée. Ligne du milieu (gauche à droite) :
           matrice de co-occurrence GLCM, carte LBP, carte d'énergie HOG. Ligne du bas :
           sous-bandes DWT (LL, LH, HL, HH).}
  \label{fig:descriptor_demo}
\end{figure}
